\section{Rancangan Algoritma CoDA-EDA}

Berdasarkan tiga prinsip rancangan sebelumnya, subbab ini menjabarkan CoDA-EDA secara konkret, meliputi representasi solusi dan model probabilistik, inisialisasi dan \textit{sampling}, pembaruan parameter, pembaruan struktur dependensi, algoritma keseluruhan, serta kompleksitas memori dan waktu. 
Rancangan ini merupakan hasil aktivitas \textit{Design and Development} yang dapat disempurnakan berdasarkan hasil \textit{Demonstration} dan \textit{Evaluation} sesuai sifat iteratif DSRM.

Ketiga unsur utama CoDA-EDA diwujudkan secara langsung dalam rancangan algoritma. 
Unsur \textit{compact} diterapkan melalui pembaruan \textit{Probability Vector} (PV) tanpa populasi eksplisit, yang membedakannya dari BOA. 
Unsur \textit{dependency-aware} diwujudkan melalui pohon Chow–Liu orde satu yang dibangun menggunakan algoritma Prim agar kebutuhan memorinya tetap $O(D)$, berbeda dari asumsi independensi \textit{compact-GA/PSO} maupun jaringan Bayes multivariat BOA. 
Unsur ketiga adalah pemisahan pembaruan parameter yang dilakukan setiap siklus dari pembaruan struktur yang hanya dipicu ketika \textit{drift} melampaui ambang $\tau$, yang membedakannya dari iBOA.

Setiap kandidat aturan ABAC direpresentasikan sebagai vektor biner
$
x=(x_1,x_2,\ldots,x_D),
$\label{first:symbol-x}
dengan $D$ menyatakan jumlah pasangan atribut–nilai dan aksi yang mungkin. 
Nilai $x_i=1$ menunjukkan bahwa pasangan ke-$i$\label{first:symbol-i} digunakan dalam aturan, sedangkan $x_i=0$ menunjukkan sebaliknya. 
Kebijakan multi-aturan dibentuk menggunakan strategi \textit{separate-and-conquer}: algoritma mencari satu aturan, mengeluarkan \textit{permit} yang telah tercakup, kemudian mengulangi proses hingga anggaran evaluasi habis atau tidak diperoleh tambahan cakupan. 
Semantik aturan mengikuti ABAC Lab, yaitu nilai dalam satu atribut bersifat disjungtif, antaratribut bersifat konjungtif, dan atribut tanpa nilai terpilih dianggap tidak dibatasi \parencite{bui_abac_2025}. 
WSC suatu aturan sama dengan jumlah pasangan atribut–nilai aktif dalam vektor tersebut.

Kualitas kandidat dievaluasi melalui fungsi \textit{fitness}
$
f(x)=\frac{TP^2}{TP+w_{fp}FP}-w_{wsc}WSC(x),
$\label{first:symbol-fx}
dengan $TP$ menyatakan jumlah \textit{permit} baru yang dicakup, $FP$ jumlah \textit{deny} yang keliru diizinkan, dan $WSC(x)$ kompleksitas aturan. 
Suku pertama ekuivalen dengan $TP$ dikalikan presisi sehingga secara simultan menghargai cakupan dan ketepatan aturan. 
Kandidat dengan $TP=0$ diberi nilai yang sangat rendah agar tidak terpilih. 
Formulasi ini menggantikan penalti linier $TP-w_{fp}FP$, yang pada pengujian awal cenderung menghasilkan aturan sangat spesifik dengan cakupan rendah. 
Keputusan ini sejalan dengan pentingnya ukuran kualitas aturan dalam mengarahkan proses \textit{separate-and-conquer} \parencite{gudys_separate_2024}.

Inisialisasi PV dibuat sadar-dimensi karena aturan ABAC yang baik umumnya relatif ringkas, sedangkan inisialisasi $p_i=0{,}5$\label{first:symbol-pi} menghasilkan sekitar separuh dari $D$ variabel aktif dan cenderung membentuk aturan terlalu spesifik. 
Oleh karena itu, setiap elemen PV diinisialisasi dengan
$
p_i=p_0=\frac{k}{D},
$\label{first:symbol-p0}
dengan $k\approx3$, sehingga jumlah atribut-nilai aktif yang diharapkan tetap relatif konstan ketika dimensi berubah. 
Pilihan ini konsisten dengan penggunaan kesederhanaan struktural sebagai salah satu ukuran kualitas kebijakan ABAC \parencite{perez-haro_abac_2024,quan_attribute-based_2023}. 
Pencarian juga dibenihkan dari satu contoh \textit{permit} yang belum tercakup untuk menghindari wilayah \textit{fitness} datar akibat rendahnya peluang aturan acak mencocokkan data, sesuai prinsip pembelajaran berbasis contoh positif pada \textit{sequential covering} \parencite{gudys_separate_2024}.

Model probabilistik CoDA-EDA terdiri atas PV dan struktur dependensi $T$\label{first:tgede}. PV,
$
PV=(p_1,p_2,\ldots,p_D),
$
merepresentasikan probabilitas marginal setiap variabel serta populasi virtual berukuran $NP$, mengikuti prinsip algoritma \textit{compact} \parencite{khalfi_metaheuristics_2023}. 
Sementara itu, $T$ merupakan pohon Chow–Liu yang mempertahankan dependensi berpasangan dengan \textit{mutual information} tertinggi \parencite{chow_approximating_1968}. 
Struktur pohon membatasi setiap variabel nonakar pada satu induk, sehingga lebih sederhana dibandingkan \textit{Bayesian network} multivariat pada BOA/iBOA \parencite{pelikan_iboa_2008}.

Pada awal pencarian, PV ditetapkan sebesar $p_i=k/D$ dan struktur $T$ belum memuat dependensi. 
Solusi elit awal dibentuk dari contoh \textit{permit} yang belum tercakup. 
Setelah struktur tersedia, kandidat baru di-\textit{sampling} menurut urutan ancestral $T$: variabel akar dibangkitkan berdasarkan probabilitas marginal, sedangkan variabel anak menggunakan probabilitas bersyarat terhadap nilai induknya.

Parameter probabilistik diperbarui pada setiap iterasi menggunakan kompetisi pemenang–pecundang seperti pada \textit{compact-GA} dan iBOA. 
Kandidat dan elit dibandingkan berdasarkan \textit{fitness}, kemudian probabilitas variabel diperbarui sebesar $\pm1/NP$ sesuai perbedaan nilai antara pemenang dan pecundang. 
Kandidat menggantikan elit apabila memiliki \textit{fitness} lebih baik. 
Karena tidak ada populasi eksplisit yang disimpan, mekanisme ini mempertahankan kebutuhan memori linear terhadap $D$.

Berbeda dengan parameter, struktur dependensi $T$ hanya diperbarui ketika terjadi perubahan distribusi yang signifikan. 
Pergeseran distribusi pasangan atribut–nilai antara jendela berjalan dan jendela referensi diukur sebagai $\Delta$\label{first:symbol-delta}. 
Jika $\Delta>\tau$, struktur Chow–Liu dipelajari ulang melalui \textit{mutual information} berpasangan dan \textit{maximum spanning tree}, sedangkan pembaruan parameter tetap berlangsung pada setiap siklus.

Pembelajaran struktur dirancang agar tidak menghilangkan sifat \textit{compact}. 
\textit{Maximum spanning tree} dibangun menggunakan algoritma Prim \parencite{cormen_introduction_2022} dengan menyimpan hanya larik sepanjang $D$, sementara \textit{mutual information} dihitung sesuai kebutuhan tanpa menyimpan matriks $D\times D$. 
Dengan demikian, kebutuhan memori tambahan tetap $O(D)$. 
Implementasi awal yang menyimpan seluruh matriks \textit{mutual information} membutuhkan sekitar 29 MB pada $D=1.907$, sehingga tidak digunakan. 
Perhitungan \textit{mutual information} juga menggunakan presisi ganda untuk menghindari hilangnya presisi numerik ketika probabilitas mendekati batas 0 atau 1.

Pemisahan pembaruan parameter dan struktur merupakan mekanisme utama CoDA-EDA: parameter beradaptasi pada setiap siklus, sedangkan struktur hanya dipelajari ulang ketika perubahan distribusi cukup besar. 
Mekanisme ini menjadi dasar pengujian H3 dan H4. 
Arsitektur keseluruhan, yang mencakup model probabilistik, jalur pembaruan parameter, dan jalur pembaruan struktur berbasis \textit{drift}, ditunjukkan pada Gambar \ref{fig:arsitektur-codaeda}, sedangkan langkah algoritmiknya dirinci pada Algoritma \ref{alg:coda-eda}.

\begin{figure}[H] 
    \centering
    \includegraphics[width=1\textwidth]{flowchart-coda-eda.png}
    \caption{Diagram Alir Arsitektur CoDA-EDA}
    \label{fig:arsitektur-codaeda}
\end{figure}

\begin{algorithm}[htbp]
\caption{CoDA-EDA (Compact Dependency-Aware EDA)}
\label{alg:coda-eda}
\begin{algorithmic}[1]
\Statex \textbf{Input}: NP (ukuran populasi virtual), $\tau$ (ambang delta pemicu pembaruan struktur), kriteria pemberhentian
\Statex \textbf{Output}: kebijakan ABAC elite ($\mathbf{x}_{\text{elite}}$)

\State Inisialisasi $\mathbf{PV} \leftarrow (0{,}5, 0{,}5, \ldots, 0{,}5)$; $T \leftarrow \emptyset$\label{first:symbol-empty} (struktur kosong)
\State $\mathbf{x}_{\text{elite}} \leftarrow \text{SAMPLE}(\mathbf{PV})$
\While{kriteria pemberhentian belum terpenuhi}
    \State $\mathbf{x}_{\text{baru}} \leftarrow \text{SAMPLE}(\mathbf{PV}, T)$ \Comment{urutan ancestral mengikuti $T$}
    \State evaluasi $f(\mathbf{x}_{\text{baru}})$ dan $f(\mathbf{x}_{\text{elite}})$
    \State $(\text{winner}, \text{loser}) \leftarrow \text{TURNAMEN}(\mathbf{x}_{\text{baru}}, \mathbf{x}_{\text{elite}})$
    \For{$i \gets 1$ \textbf{to} $D$}
        \State $\mathbf{PV}[i] \leftarrow \mathbf{PV}[i] \pm \dfrac{1}{NP}$ berdasarkan $\text{winner}[i]$, $\text{loser}[i]$
    \EndFor
    \If{$f(\mathbf{x}_{\text{baru}})$ lebih baik dari $f(\mathbf{x}_{\text{elite}})$}
        \State $\mathbf{x}_{\text{elite}} \leftarrow \mathbf{x}_{\text{baru}}$
    \EndIf
    \State $\Delta \leftarrow \text{UKUR\_DRIFT}(\text{jendela akses berjalan}, \text{jendela referensi})$
    \If{$\Delta > \tau$}
        \State $T \leftarrow \text{CHOW-LIU}(\text{mutual information berpasangan}, \text{pohon rentang maksimum})$
        \State perbarui jendela referensi
    \EndIf
\EndWhile
\State \Return{$\mathbf{x}_{\text{elite}}$}
\end{algorithmic}
\end{algorithm}\label{first:symbol-empty}

Algoritma \ref{alg:coda-eda} merangkum keseluruhan mekanisme CoDA-EDA. 
Baris 1–2 memuat inisialisasi, baris 4–12 menjalankan \textit{sampling} dan pembaruan parameter pada setiap iterasi, sedangkan baris 13–17 menangani deteksi \textit{drift} dan pembaruan struktur secara bersyarat.

Analisis kompleksitas berikut diturunkan berdasarkan rancangan algoritma \textit{(by construction)} dan selanjutnya divalidasi secara empiris melalui eksperimen. 
Penelitian ini tidak mencakup pembuktian konvergensi formal secara mendalam.

Kompleksitas memori CoDA-EDA adalah $O(D)$. 
PV membutuhkan $O(D)$, sedangkan struktur pohon $T$ juga membutuhkan $O(D)$ karena setiap variabel hanya menyimpan satu referensi induk dan tabel probabilitas bersyarat berukuran konstan. 
Dengan demikian, kebutuhan memori tetap linear terhadap dimensi dan tidak bergantung pada ukuran populasi virtual (NP). 
Karakteristik ini sebanding dengan algoritma \textit{compact} standar, tetapi lebih terbatas dibandingkan iBOA yang dapat memerlukan hingga $O(n^k)$ tabel probabilitas ketika orde dependensi $k$ tidak dibatasi.

Untuk analisis waktu, $D$ menyatakan dimensi ruang solusi, $n$\label{first:symbolnn} jumlah entri \textit{log} per jendela, $E$ anggaran evaluasi \textit{fitness} per jendela, $W$ jumlah jendela, $\delta$\label{first:symbol-deltagede} proporsi jendela yang memicu pembaruan struktur, dan $k$ rata-rata jumlah kondisi aktif dalam aturan. 
Karena aturan ABAC bersifat jarang, $k$ diharapkan jauh lebih kecil daripada $D$.

Satu evaluasi kandidat terdiri atas tiga tahap. 
Pertama, kandidat dibangkitkan melalui ancestral \textit{sampling} pada pohon $T$ dengan biaya $O(D)$. 
Kedua, \textit{fitness} dihitung dengan memindai $n$ entri \textit{log} untuk memperoleh $TP$ dan $FP$. 
Pada implementasi umum, biaya terburuknya adalah $O(n \cdot D)$, sedangkan dengan hanya memeriksa $k$ kondisi aktif dapat mendekati $O(n \cdot k)$. 
Ketiga, PV diperbarui dalam $O(D)$. 
Dengan demikian, kompleksitas terburuk satu evaluasi adalah $O(n \cdot D)$, sehingga $E$ evaluasi membutuhkan
$
O(E \cdot n \cdot D)
$
per jendela.

Jalur pembaruan struktur diawali dengan deteksi \textit{drift}. 
Distribusi pasangan atribut–nilai pada jendela berjalan dihitung dan dibandingkan dengan jendela referensi menggunakan JSD, dengan biaya terburuk $O(n \cdot D)$. 
Jika $\Delta>\tau$, struktur Chow–Liu dipelajari ulang. 
Perhitungan \textit{mutual information} untuk seluruh pasangan variabel membutuhkan $O(n \cdot D^2)$, sedangkan pembentukan \textit{maximum spanning tree} dengan algoritma Prim pada graf lengkap membutuhkan $O(D^2)$. 
Oleh karena itu, biaya pembelajaran ulang struktur didominasi oleh
$
O(n \cdot D^2).
$

Karena pembelajaran struktur hanya terjadi pada proporsi $\delta$ dari seluruh jendela, biaya total untuk $W$ jendela dapat dinyatakan secara ringkas sebagai
$
O\left(W \cdot E \cdot n \cdot D + W \cdot n \cdot D + \delta W \cdot n \cdot D^2\right),
$
atau, dengan mengabaikan suku yang terdominasi,
$
O\left(W \cdot E \cdot n \cdot D + \delta W \cdot n \cdot D^2\right).
$

Dengan demikian, jalur parameter tetap linear terhadap $D$ untuk pembaruan model dan didominasi oleh biaya evaluasi \textit{fitness}, sedangkan biaya kuadratik terhadap dimensi hanya muncul ketika struktur dependensi perlu dipelajari ulang. 
Pemisahan ini merupakan tujuan utama mekanisme berbasis \textit{drift}: menghindari biaya $O(n \cdot D^2)$ pada jendela yang distribusinya masih stabil.

\begin{longtable}{|>{\raggedright\arraybackslash}p{2.3cm}|>{\raggedright\arraybackslash}p{1.8cm}|>{\raggedright\arraybackslash}p{1.6cm}|>{\raggedright\arraybackslash}p{2.2cm}|>{\raggedright\arraybackslash}p{3cm}|>{\raggedright\arraybackslash}p{1.5cm}|}
\caption{Perbandingan kompleksitas waktu dan memori per siklus/generasi}
\label{tab:perbandingan-kompleksitas} \\
\hline
\textbf{Algoritma} & \textbf{Sampling 1 kandidat} & \textbf{Evaluasi \textit{fitness}} & \textbf{Pembaruan model per siklus} & \textbf{Pembaruan struktur} & \textbf{Memori} \\
\hline
\endfirsthead

\multicolumn{6}{c}{{\bfseries \tablename\ \thetable{} -- lanjutan dari halaman sebelumnya}} \\
\hline
\textbf{Algoritma} & \textbf{Sampling 1 kandidat} & \textbf{Evaluasi \textit{fitness}} & \textbf{Pembaruan model per siklus} & \textbf{Pembaruan struktur} & \textbf{Memori} \\
\hline
\endhead

\hline
\multicolumn{6}{r}{Lanjut ke halaman berikutnya} \\
\endfoot

\hline
\endlastfoot

% ==================== DATA BARIS ====================
compact-GA / compact-PSO 
  & $O(D)$ 
  & $O(n \cdot D)$ 
  & $O(D)$ (PV) 
  & --- (tanpa struktur) 
  & $O(D)$ \\
\hline

CoDA-EDA 
  & $O(D)$ (ancestral pd pohon) 
  & $O(n \cdot D)$ 
  & $O(D)$ (PV) 
  & $O(n \cdot D^2)$ --- hanya saat \textit{drift} 
  & $O(D)$ \\
\hline

iBOA 
  & $O(D)$ 
  & $O(n \cdot D)$ 
  & $O(D)$ 
  & $O(D^2)$ tiap langkah (kontinu) 
  & $O(D)$ \\
\hline

BOA 
  & $O(NP \cdot D)$ 
  & $O(NP \cdot n \cdot D)$ 
  & --- 
  & $O(NP \cdot D^2)$ bangun ulang tiap generasi 
  & $O(NP \cdot D)$ \\
\hline

ACO / UBK 
  & $O(NP \cdot D)$ (konstruksi) 
  & $O(NP \cdot n \cdot D)$ 
  & $O(NP \cdot D)$ (feromon) 
  & --- (tak inkremental) 
  & $O(NP \cdot D)$ \\
\hline

\end{longtable}

Terbaca bahwa jalur per siklus CoDA-EDA satu orde dengan keluarga \textit{compact} $O(n \cdot D)$, sedangkan operasi struktur $O(n \cdot D^2)$-nya hanya dibayar saat \textit{drift}.
Bandingkan dengan iBOA yang membayar biaya struktur $O(D^2)$ pada setiap langkah (kontinu) dan dengan BOA/ACO/UBK yang seluruh biayanya diperbesar faktor populasi NP.
Perbedaan inilah yang membuat CoDA-EDA lebih efisien pada pembaruan berkala sekaligus tetap hemat memori.